\documentclass{report}
\usepackage{amsmath,amssymb}
\setlength{\parindent}{0mm}
\setlength{\parskip}{1em}
\begin{document}
\begin{center}
\rule{6in}{1pt} \
{\large Matthias Bolten \\
{\bf Local Fourier analysis of pattern structured operators}}

Bergische Universit�t Wuppertal \\ Fachbereich C - Mathematik und Naturwissenschaften \\ Gau�stra�e 20 \\ D-42119 Wuppertal \\ Germany
\\
{\tt bolten@math.uni-wuppertal.de}\\
Karsten Kahl\\
Hannah Rittich\end{center}

Multigrid methods are used to compute the solution $u$
of the system of equations
\[
L u = f\,,
\]
where $L$ is typically a discretization of a partial differential equation
(PDE) and $f$ a corresponding given right hand side.
Local Fourier Analysis (LFA) is well known to provide quantitative estimates for the
speed of convergence of multigrid methods by analyzing the involved
operators in the frequency domain.

For the initial formulation of LFA it was crucial to assume that
all involved operators have constant coefficients. In many applications the
coefficients vary continuously in space. Thus, asymptotically or if the
grid is fine enough the
discrete operator $L$ will only vary slightly between neighboring grid points
and hence can be well approximated by an operator with {\em locally} constant
coefficients. Thus constant coefficient often are a reasonable assumption.

However, when analyzing problems with jumping coefficients or when the
variance of the coefficient is large this assumption is too restrictive.
A reason for this is that interpolation and restriction
operators typically act differently on variables that have a coarse grid
representative and those who do not have one. Further, the analysis of
multigrid methods itself require an extension of this straight forward
approach, e.g. pattern
relaxation schemes like the red-black Gau�-Seidel where red points of
the grid are treated differently than black ones.

It is possible to analyze the latter when
allowing for interaction of certain frequencies. Furthermore, it turns out that when
allowing for more frequencies to interact operators given by
increasingly complex patterns can be analyzed.

In this talk a general framework for analyzing pattern
structured operators, i.e., operators whose action is invariant under certain
shifts of the input function, will be introduced and different
applications will be discussed.


\end{document}
